\chapter{Összegzés}
\label{ch:sum}

A szakdolgozat a Rush Hour logikai játék keretein belül valósította meg a keresőalgoritmusokat, heurisztikákat, a megjelenítést és az interakciót. Ezen kívül számos különböző problémára, játékra lehet érdemes hasonló elemzéseket elvégezni.
Az állapottérbeli keresésre 3 alapvető algoritmus került implementálásra, amikkel jól ábrázolhatóak a különböző megközelítési egy problémának. Arra is alkalmas a rendszer, hogy egy saját algoritmust látványosan tudjunk futtatni, és elemezzük annak viselkedését, ellenőrizzük az elvárt helyességét. 

\paragraph{Bővítési lehetőségek.} Számos területen lehet még tovább fejlesztheti az alkalmazást:
\begin{compactitem}
\item A Rush Hour-ön túlmutató további rejtvénytípusok támogatása a keretrendszer általánosítása érdekében.
\item Az A$^*$ keresés implementálása (a lezárt állapotok újranyitásával) az A$^C$ algoritmussal való összehasonlítás céljából.
\item Párhuzamos algoritmusfuttatás a különböző konfigurációk egymás melletti összehasonlításához.
\item GPU-gyorsított erőszámítás a több ezer csúcsból álló gráfok kezeléséhez.
\item Algoritmusfutások rögzítése és visszajátszása oktatási prezentációkhoz.
\end{compactitem}